\documentclass[11pt,letterpaper]{article}
\usepackage[utf8]{inputenc}
\usepackage[T1]{fontenc}
\usepackage{geometry}
\geometry{margin=1in}
\usepackage{hyperref}

\pagestyle{empty}

\begin{document}

\today

\vspace{1cm}

Dear Editor,

\vspace{0.5cm}

We submit for your consideration the perspective article entitled ``\textbf{Memory as a Biological Property: How a Single Number from Time-Series Econometrics Reveals Hierarchy in Gene Expression}'' for publication in \textit{Cell Systems}.

\vspace{0.3cm}

\textbf{Summary.} We argue that temporal persistence---how much a gene's past determines its future---is a fundamental axis of biological variation, distinct from expression level, network connectivity, or epigenetic state. The eigenvalue modulus of a second-order autoregressive model provides a direct measurement of this property using mathematics that has been refined across physics, engineering, economics, and climate science over two centuries. Evidence from 37 datasets across 4 species demonstrates that this single metric blindly recovers known circadian hierarchy, is partially independent of chromatin state and network connectivity, and converges with four independent research programs.

\vspace{0.3cm}

\textbf{Significance.} This perspective connects gene expression dynamics to the broader scientific tradition of eigenvalue analysis, showing that the same mathematical object that determines energy levels in quantum mechanics, natural frequencies in structural engineering, and economic persistence timescales also measures biological memory in gene expression. We present evidence that eigenvalue is partially independent of network connectivity ($\rho = -0.29$) and nearly independent of chromatin state ($\rho = 0.08$), confirming it captures a distinct biological dimension. All results and interactive tools are available at \url{https://par2-discovery-engine.replit.app}.

\vspace{0.3cm}

\textbf{Companion paper.} The quantitative methods and comprehensive robustness validation underlying this perspective are detailed in the companion paper ``AR(2) Eigenvalue Modulus as a Measure of Temporal Persistence in Gene Expression'' (Whiteside, 2026), which presents the full eleven-analysis robustness suite, five ODE model validations, and cross-species replication across 35 tissue-level time series.

\vspace{0.3cm}

\textbf{Honest limitations.} We explicitly acknowledge that this perspective piece argues for a framework rather than proving causation; that the AR(2) model assumes linearity; that the aging and disease predictions remain speculative; and that evidence is predominantly from circadian contexts. These limitations are discussed in detail in the manuscript.

\vspace{0.3cm}

\textbf{Broad appeal.} This article bridges computational biology, time-series statistics, dynamical systems theory, and circadian biology. It should interest researchers across these fields who may not be aware that their mathematical tools apply directly to gene expression analysis.

\vspace{0.3cm}

This manuscript has not been submitted elsewhere and all work is original.

\vspace{0.5cm}

Sincerely,

\vspace{0.5cm}

Michael Whiteside \\
Independent Researcher, United Kingdom \\
mickwh@msn.com \\
ORCID: 0009-0000-0643-5791

\end{document}
